\section{Design Thinking \& Storyboard}

This section documents the design thinking process used by the team to
arrive at the final Smart Cart concept. Design thinking was applied across
five stages: Empathize, Define, Ideate, Prototype, and Test. Each stage
shaped key decisions about the system's hardware configuration, user
interaction model, and safety approach.

\subsection{Empathize --- Understanding the Problem}

The team began by observing shopping behavior in a large retail mall
environment. Key pain points identified through observation and informal
user interviews included:
\begin{itemize}
  \item Elderly shoppers struggled to push heavily loaded carts, especially on slight inclines near entrances.
  \item Parents with young children frequently lose control of carts when distracted.
  \item Narrow aisles in peak hours caused collisions between carts and other shoppers, particularly at ankle and heel height.
  \item Cart return was often neglected, causing trolleys to roll freely in car parks.
\end{itemize}

These direct observations motivated the core design goal: a cart that
removes physical burden from the user while actively preventing collisions
in dynamic environments.

\subsection{Define --- Problem Statement}

Based on the empathy stage, the team defined the central problem as:

\begin{quote}
\textit{``Shoppers in crowded retail environments need a smarter cart
that follows them autonomously, detects obstacles, and adapts to load ---
so they can shop safely without physical strain.''}
\end{quote}

\subsection{Ideate --- Design Concepts Considered}

The team generated three competing design concepts before converging on
the final solution:

\begin{table}[!ht]
\caption{Design Concepts Considered During Ideation (Author developed)}
\label{tab:ideate}
\centering
\footnotesize
\begin{tabular}{p{1.2cm} p{2.3cm} p{3.0cm}}
\toprule
\textbf{Concept} & \textbf{Description} & \textbf{Why Rejected / Adopted} \\
\midrule
Concept A & GPS-guided outdoor cart with motorized wheels & GPS unreliable indoors; rejected for indoor-only scope \\
Concept B & Camera-only follow system (no UWB) & Camera occlusion in crowds caused unreliable tracking; rejected \\
Concept C (Chosen) & UWB-based follow-me with omni-wheels, ultrasonic + IR sensors & Most reliable indoors; handles occlusion; low cost; ADOPTED \\
\bottomrule
\end{tabular}
\end{table}